% !TeX root = ../main.tex

\chapter{Differentiation}

% =================================================

Let

\[
\phi_n(x)=e^{\lambda x^n}.
\]

Differentiating gives

\[
\phi_n'(x)
=
n\lambda x^{n-1}e^{\lambda x^n}.
\]

Therefore

\[
\boxed{
\phi_n'(x)
=
n\lambda x^{n-1}\phi_n(x).
}
\]

The original family

\[
\{
\phi_n
\}
\]

is therefore not closed under differentiation because a polynomial
factor appears.

However, this suggests enlarging the family.

Consider functions of the form

\[
x^k\phi_n(x).
\]

Differentiate:

\[
\frac{d}{dx}
\left(
x^k\phi_n(x)
\right)
=
kx^{k-1}\phi_n(x)
+
x^k\phi_n'(x).
\]

Therefore

\[
\boxed{
\frac{d}{dx}
\left(
x^k\phi_n(x)
\right)
=
kx^{k-1}\phi_n(x)
+
n\lambda x^{k+n-1}\phi_n(x).
}
\]

Both terms remain in the enlarged family

\[
\boxed{
\mathcal{D}
=
\{
x^k e^{\lambda x^n}
:
k\geq0,\ n\geq1
\}.
}
\]

Thus this larger family \emph{is} closed under differentiation.

% =================================================
\section{Connection with Differential Equations}
% =================================================

Each block satisfies a first-order differential equation.

From

\[
\phi_n'(x)
=
n\lambda x^{n-1}\phi_n(x),
\]

we obtain

\[
\boxed{
\phi_n'(x)
-
n\lambda x^{n-1}\phi_n(x)
=
0.
}
\]

Define the differential operator

\[
L_n
=
\frac{d}{dx}
-
n\lambda x^{n-1}.
\]

Then

\[
\boxed{
L_n\phi_n=0.
}
\]

Thus every basic block is the solution of a simple first-order linear
differential equation.

This gives another possible direction:

\begin{question}

Can properties of a monomial exponential series be described using
the collection of differential operators

\[
L_1,L_2,L_3,\ldots?
\]

\end{question}

% =================================================
\section{Multiplication of Blocks}
% =================================================

Let us also investigate multiplication.

Take

\[
\phi_m(x)=e^{\lambda x^m}
\]

and

\[
\phi_n(x)=e^{\lambda x^n}.
\]

Then

\[
\phi_m(x)\phi_n(x)
=
e^{\lambda x^m}
e^{\lambda x^n}.
\]

Therefore

\[
\boxed{
\phi_m(x)\phi_n(x)
=
e^{\lambda(x^m+x^n)}.
}
\]

In general this is not another basic block.

Thus

\[
\{\phi_n\}
\]

is not closed under multiplication.

However, multiplication naturally leads to the larger family

\[
\boxed{
e^{\lambda P(x)},
}
\]

where \(P(x)\) is a polynomial.

For example,

\[
\phi_2(x)\phi_5(x)
=
e^{\lambda(x^2+x^5)}.
\]

This suggests that the monomial exponential blocks sit naturally
inside a larger theory of exponentials of polynomials.

% =================================================
